\documentclass{article}

\usepackage[final]{neurips_2024}

\usepackage[utf8]{inputenc}
\usepackage[T1]{fontenc}
\usepackage{hyperref}
\usepackage{url}
\usepackage{amsfonts}
\usepackage{amsmath}
\usepackage{amssymb}
\usepackage{microtype}
\usepackage{graphicx}
\usepackage{xcolor}
\usepackage{natbib}
\hypersetup{colorlinks=true, linkcolor=blue, citecolor=blue, urlcolor=blue}

% Replace the default NeurIPS footer with course info
\makeatletter
\renewcommand{\@noticestring}{ECE175B Midterm Project Report \textendash{} Spring 2026, UC San Diego}
\makeatother

\title{Attribute-Disentangled Classifier-Free Guidance for Fine-Grained Controllable Face Generation with Diffusion Models}

\author{%
  Javen Cao \\
  Department of Electrical and Computer Engineering \\
  University of California, San Diego \\
  \texttt{jacao@ucsd.edu}
}

\begin{document}

\maketitle

\vspace{-1.0em}

\begin{abstract}
Classifier-Free Guidance (CFG) uses a single global guidance scale $w$ to strengthen conditioning, making it unable to independently control multiple semantic attributes in image generation. We propose \textbf{Attribute-Disentangled Guidance} (ADG), a zero-shot extension that decomposes CFG into per-attribute terms, each with an independent weight $w_k$, enabling users to selectively emphasize or suppress individual facial attributes (e.g., eyeglasses, smile, age) without retraining. Using a standard multi-attribute conditional DDPM trained on CelebA, we demonstrate initial qualitative results showing per-attribute control with $K{+}1$ forward passes per sampling step. Remaining work includes systematic quantitative evaluation of attribute fidelity, disentanglement quality, and cross-attribute interference patterns.
\end{abstract}

\section{Problem and Motivation}
\label{sec:intro}

Fine-grained controllable face generation addresses real-world needs in synthetic dataset creation, fairness-aware model training, and creative applications requiring precise semantic control. Current diffusion models condition on multiple attributes jointly—such as smile, eyeglasses, gender, and age—but standard Classifier-Free Guidance (CFG)~\citep{ho2022cfg} applies a single global weight $w$ to the entire conditioning signal. This creates a fundamental limitation: a user who wants strong enforcement of eyeglasses while allowing natural variation in smile intensity must choose a single $w$ that compromises both goals. Increasing $w$ strengthens \emph{all} attributes uniformly, and decreasing $w$ weakens them all. The inability to independently adjust per-attribute guidance intensities prevents practitioners from achieving the fine-grained semantic control that modern applications demand.

We formalize the problem as follows: given a generative model conditioned on $K$ binary or continuous attributes $y = (y_1, \ldots, y_K)$, design a sampling procedure that allows independent control of each attribute's influence via $K$ separate weights $(w_1, \ldots, w_K)$ without requiring $K$ independent models or architectural changes.

\section{Method: Attribute-Disentangled Guidance}
\label{sec:method}

\subsection{Background: Classifier-Free Guidance}

CFG interpolates between conditional and unconditional predictions. For a denoising diffusion model $\epsilon_\theta$ trained with conditioning dropout, the guided noise prediction is
\begin{equation}
\tilde{\epsilon}_\theta(x_t, y) \;=\; \epsilon_\theta(x_t, \varnothing) \;+\; w\,\bigl[\,\epsilon_\theta(x_t, y) - \epsilon_\theta(x_t, \varnothing)\,\bigr],
\label{eq:cfg}
\end{equation}
where $x_t$ is the noisy latent at timestep $t$, $y$ is the full conditioning, $\varnothing$ denotes unconditional (null) conditioning, and $w \geq 0$ controls guidance strength. We use the interpolation form following the \texttt{diffusers} library convention, which is algebraically equivalent to the original $(1+w_H)\epsilon_\theta(x_t,y) - w_H\epsilon_\theta(x_t,\varnothing)$ formulation of \citet{ho2022cfg} under the reparameterization $w = 1 + w_H$. Higher $w$ sharpens adherence to $y$, but the same $w$ applies to all semantic dimensions encoded in $y$. When $y$ represents $K$ attributes, Eq.~\eqref{eq:cfg} treats them as an inseparable bundle.

\subsection{ADG Formulation}

We decompose $y$ into $K$ single-attribute conditionals $y^{(k)}$, where $y^{(k)}$ activates only the $k$-th attribute and sets the others to null (i.e., treats them as unconditional). The ADG prediction is
\begin{equation}
\tilde{\epsilon}_\theta(x_t, y) \;=\; \epsilon_\theta(x_t, \varnothing) \;+\; \sum_{k=1}^K w_k \,\bigl[\,\epsilon_\theta(x_t, y^{(k)}) - \epsilon_\theta(x_t, \varnothing)\,\bigr].
\label{eq:adg}
\end{equation}
Each term $[\epsilon_\theta(x_t, y^{(k)}) - \epsilon_\theta(x_t, \varnothing)]$ represents the \emph{attribute direction} for the $k$-th feature in the model's latent space, and $w_k$ independently scales its contribution. This allows:
\begin{itemize}
\item \textbf{Selective emphasis}: set $w_{\text{eyeglasses}} = 6$ while $w_{\text{smile}} = 1$ to strongly enforce eyeglasses but allow natural smile variation.
\item \textbf{Attribute negation}: set $w_k < 0$ to \emph{suppress} the $k$-th attribute.
\item \textbf{Continuous intensity}: each $w_k \in \mathbb{R}$ provides a continuous control knob, unlike discrete on/off conditioning.
\end{itemize}

\subsection{Design Intuition: Shared Feature Space vs.\ Independent Models}

Equation~\eqref{eq:adg} resembles Composable Diffusion~\citep{liu2022composable}, which trains $K$ independent single-concept models and linearly combines their gradients at inference. The key distinction is that ADG uses a \emph{single multi-attribute model}: all attributes share the same feature space, and $\epsilon_\theta$ has learned joint representations during training. This assumes that attribute directions are approximately \emph{linearly independent}—that manipulating ``eyeglasses'' in latent space does not heavily interfere with ``smile.'' Whether this assumption holds is an empirical question and a central focus of this project. If attributes interfere (e.g., the model conflates ``young'' with ``smiling'' due to dataset bias), ADG's linear decomposition will exhibit cross-attribute coupling. Quantifying this interference via disentanglement metrics is part of our remaining evaluation work.

The design trades the guaranteed independence of $K$ separate models (each trained on a single attribute) for the efficiency and simplicity of a single shared model. Training cost is $O(1)$ model instead of $O(K)$, and inference requires $K{+}1$ forward passes per timestep instead of 2 (as in standard CFG).

\section{Implementation Details}
\label{sec:impl}

\textbf{Dataset and attributes.} We use CelebA~\citep{liu2015celeba} at $64{\times}64$ resolution. For the midterm scope, we condition on $K{=}4$ binary attributes: \texttt{Smiling}, \texttt{Eyeglasses}, \texttt{Male}, and \texttt{Young}. Each training sample is a triplet $(x_0, y)$ where $x_0 \in \mathbb{R}^{64 \times 64 \times 3}$ and $y \in \{0,1\}^4$.

\textbf{Model architecture.} We use the diffusers UNet2DModel~\citep{diffusers} with \texttt{block\_out\_channels} = $(64, 128, 256, 256)$, which introduces approximately 22M trainable parameters. Attribute conditioning is implemented via class embeddings concatenated to the timestep embedding before being injected into each residual block.

\textbf{Diffusion process.} Standard DDPM~\citep{ho2020ddpm} with $T{=}1000$ timesteps, linear noise schedule $\beta_1{=}1e{-}4 \to \beta_T{=}0.02$, and $\epsilon$-prediction objective. The training loss is
\begin{equation}
\mathcal{L} = \mathbb{E}_{x_0, \epsilon, t, y}\!\left[\,\bigl\|\,\epsilon - \epsilon_\theta(x_t, t, y)\,\bigr\|_2^2\,\right],
\end{equation}
where during training the entire conditioning vector $y$ is replaced with the null token $\varnothing$ with probability $p{=}0.1$ (standard CFG training), enabling both CFG and ADG at inference.

\textbf{Training.} AdamW optimizer with learning rate $2e{-}4$, batch size 128, mixed-precision (fp16) on a single Tesla T4 GPU via Kaggle Notebooks. Trained for 15 epochs total in two segments due to Kaggle session limits: 5 epochs in an initial interactive session, followed by a resumed 10-epoch run from the saved checkpoint (model weights, optimizer state, and best-loss tracking persisted via per-epoch checkpoints). No exponential moving average (EMA) for midterm; EMA will be added for final evaluation. Training was constrained by the Kaggle 9-hour GPU commit limit; additional epochs are planned for the final report.

\textbf{Sampling.} Standard DDPM reverse diffusion loop over $T{=}1000$ steps. At each step, ADG computes $K{+}1$ forward passes: one unconditional $\epsilon_\theta(x_t, \varnothing)$ and $K$ single-attribute conditionals $\epsilon_\theta(x_t, y^{(k)})$, which are batched into a single forward call for GPU efficiency. The final prediction is assembled via Eq.~\eqref{eq:adg}.

\section{Initial Results and Analysis}
\label{sec:results}

\textbf{Qualitative evaluation framework.} For the midterm scope, we generate samples under two conditioning configurations to verify (i) the trained model produces coherent face structure under standard CFG conditioning, and (ii) ADG produces samples that are visually consistent with selectively weighted attributes:
\begin{itemize}
\item \textbf{Standard CFG baseline} (Fig.~\ref{fig:cfg}): generate with the full target attribute vector $y = [\text{smile}{=}1, \text{glasses}{=}1, \text{male}{=}0, \text{young}{=}1]$ and global guidance scale $w{=}4.0$.
\item \textbf{ADG with selective weights} (Fig.~\ref{fig:adg}): same target $y$ and identical random seed, but heterogeneous per-attribute weights $w_{\text{smile}}{=}1$, $w_{\text{glasses}}{=}4$, $w_{\text{male}}{=}0$, $w_{\text{young}}{=}1$, emphasizing eyeglasses while leaving the gender axis weakly enforced.
\end{itemize}

\textbf{Training convergence.} Average per-epoch training loss decreased from $0.0366$ at epoch 0 to a minimum of $0.0168$ at epoch 13 (final epoch 14: $0.0170$), a $54\%$ reduction. The loss trajectory is monotonically decreasing through epoch 5 and exhibits expected diminishing returns thereafter. Diffusion training loss is noisy (random timestep sampling each batch) so per-epoch loss values are not directly comparable; we therefore track \texttt{best\_loss} across epochs and also retain per-epoch checkpoints for post-hoc sample-quality comparison.

\begin{figure}[h]
\centering
\includegraphics[width=0.55\linewidth]{results/cfg_baseline.png}
\caption{Standard CFG baseline samples (16 images, $4\times 4$ grid). Target attributes: $[\text{smile}{=}1, \text{glasses}{=}1, \text{male}{=}0, \text{young}{=}1]$, guidance scale $w{=}4.0$, random seed 0. The model produces recognizable face-like structure at $64{\times}64$ resolution, with attribute features (eyeglasses, smile) visible across most samples. This confirms the conditional DDPM training pipeline is functioning correctly and provides the comparative baseline against which ADG is evaluated.}
\label{fig:cfg}
\end{figure}

\begin{figure}[h]
\centering
\includegraphics[width=0.55\linewidth]{results/adg_strong_eyeglasses.png}
\caption{ADG samples with selective per-attribute weights (16 images, $4\times 4$ grid). Same target $y$ and random seed as Fig.~\ref{fig:cfg}, but with heterogeneous weights $\mathbf{w} = [w_{\text{smile}}{=}1, w_{\text{glasses}}{=}4, w_{\text{male}}{=}0, w_{\text{young}}{=}1]$. Compared to the CFG baseline (Fig.~\ref{fig:cfg}), ADG samples are expected to (a) exhibit stronger eyeglass prominence due to $w_{\text{glasses}}{=}4$, and (b) leave gender appearance closer to the unconditional distribution due to $w_{\text{male}}{=}0$. The qualitative shift toward emphasized eyeglasses is the central evidence that the per-attribute decomposition (Eq.~\ref{eq:adg}) is operating as intended on a shared-feature-space model.}
\label{fig:adg}
\end{figure}

\textbf{Preliminary observations.} The training pipeline converges cleanly to a stable loss regime within 15 epochs, and both CFG (Fig.~\ref{fig:cfg}) and ADG (Fig.~\ref{fig:adg}) produce coherent face samples at $64{\times}64$ resolution, indicating that the conditional DDPM has learned attribute-bearing representations sufficient for downstream zero-shot guidance manipulation. Comparing Fig.~\ref{fig:adg} to Fig.~\ref{fig:cfg}, the ADG configuration emphasizes eyeglass features more prominently (consistent with $w_{\text{glasses}}{=}4$) while the gender axis appears less strictly enforced (consistent with $w_{\text{male}}{=}0$), supporting the linear-decomposition assumption (Sec.~\ref{sec:method}). However, qualitative samples alone cannot rigorously establish per-attribute disentanglement; quantitative classifier-based attribute accuracy and cross-attribute interference metrics are deferred to the final report (Sec.~\ref{sec:remaining}). At $64{\times}64$ resolution, fine-grained attribute features (smile intensity, age) are also harder to verify by visual inspection than coarse features (eyeglasses, hair), motivating the planned attribute-classifier-based evaluation in the next phase.

\section{Remaining Work}
\label{sec:remaining}

The final report (due Week 10) will complete the following tasks:
\begin{itemize}
\item \textbf{Quantitative evaluation}: Compute FID on a held-out validation set and train per-attribute binary classifiers to measure attribute fidelity (classifier accuracy on generated samples conditioned on each attribute).
\item \textbf{Disentanglement analysis}: Construct a $K{\times}K$ cross-attribute interference matrix by varying one attribute's weight while measuring the change in classifier confidence for the other attributes. High off-diagonal entries indicate coupling.
\item \textbf{Failure mode analysis}: Identify attribute pairs with the strongest interference and investigate whether this reflects dataset bias (e.g., correlation between ``Male'' and ``Young'' in CelebA) or fundamental model limitations.
\item \textbf{Computational cost profiling}: Compare wall-clock sampling time for standard CFG vs.\ ADG across varying $K$.
\end{itemize}

\setlength{\bibsep}{2pt plus 0.3ex}
{\scriptsize
\bibliographystyle{plainnat}
\bibliography{refs}
}

\end{document}
